\section*{Bab \ref{ch:g12:destruction-of-mountains} --- Hancurnya Pegunungan}

\begin{solution}{exo:g12:destruction-of-mountains:1}
(a) \omterm{def:g12:destruction-of-mountains:denudation}{laju denudasi}; (b) \omterm{def:g12:destruction-of-mountains:load}{muatan dasar}; (c) \omterm{def:g12:destruction-of-mountains:load}{muatan terlarut};
(d) \omterm{def:g12:destruction-of-mountains:diagenesis}{diagenesis}; (e) \omterm{def:g12:destruction-of-mountains:rebound}{pantulan isostatik}.
\end{solution}

\begin{solution}{exo:g12:destruction-of-mountains:2}
(1) Konglomerat. (2) Batu pasir. (3) Serpih. (4) Breksi.
\end{solution}

\begin{solution}{exo:g12:destruction-of-mountains:3}
\qty{1}{m/Ma} adalah \qty{1000}{mm} dalam \num{1000000} tahun, yaitu
\qty{0.001}{mm} setahun. (1) \qty{0.054}{mm} setahun. (2) \qty{0.218}{mm}
setahun. (3) $\qty{10}{mm}$ setahun adalah \qty{10}{m} dalam 1000 tahun, yaitu
\qty{10000}{m/Ma}.
\end{solution}

\begin{solution}{exo:g12:destruction-of-mountains:4}
Muatan padat adalah \qty{85}{\percent} dari seluruhnya, sehingga seluruhnya $60/0.85
\approx 70.6$ juta ton, dan \omterm{def:g12:destruction-of-mountains:load}{muatan terlarutnya} sekitar $70.6 - 60 \approx
10.6$ juta ton setahun.
\end{solution}

\begin{solution}{exo:g12:destruction-of-mountains:5}
Pelapukan granit, pengangkutan oleh sungai, pengendapan di sebuah cekungan,
penimbunan di bawah lapisan-lapisan yang lebih muda, sementasi.
\end{solution}

\begin{solution}{exo:g12:destruction-of-mountains:6}
Pada \qty{30}{cm/s} butir pasir terangkat (di atas \qty{20}{cm/s}); butir
lanau tetap melayang. Pada \qty{5}{cm/s} butir pasir mengendap (di bawah
\qty{10}{cm/s}), tetapi butir lanau masih tetap melayang, karena
\qty{0.1}{cm/s} sudah cukup untuk membuatnya tetap melayang.
\end{solution}

\begin{solution}{exo:g12:destruction-of-mountains:7}
Muatan seluruhnya $4/0.85 \approx 4.71$ juta ton, $\num{4.71e9}/2830 \approx
\qty{1.66e6}{m^3}$ batuan setahun. Dibagi $\qty{5e10}{m^2}$: $\num{3.3e-5}$ m
setahun, sekitar \qty{33}{m/Ma}, di bawah median \qty{54}{m/Ma}.
\end{solution}

\begin{solution}{exo:g12:destruction-of-mountains:8}
(1) $1 - 2700/3300 \approx 0.182$. (2) Sekitar \qty{182}{m}. (3) Sedikit lebih
banyak daripada \qty{163}{m} dalam bab ini: makin ringan kerak dibandingkan
dengan mantel, makin sedikit kolom itu naik.
\end{solution}

\begin{solution}{exo:g12:destruction-of-mountains:9}
(1) Batuan setebal \qty{8}{km}. (2) $8000 \times 0.163 \approx \qty{1300}{m}$.
(3) $8000/200 = 40$ juta tahun.
\end{solution}

\begin{solution}{exo:g12:destruction-of-mountains:10}
(1) $794/43 \approx 18$ juta tahun. (2) Permukaan hanya turun
$0.163$ kali tebal yang terkikis, sehingga batuan yang harus dihilangkan $794/0.163
\approx \qty{4900}{m}$, yang memerlukan waktu sekitar $4900/43 \approx 110$ juta
tahun. (3) Benua tidak hanya terkikis: pembentukan pegunungan dan magmatisme
busur terus mengangkatnya dan menambahkan kerak baru.
\end{solution}

\begin{solution}{exo:g12:destruction-of-mountains:11}
(1) $\num{1e14}/2830 \approx \qty{3.5e10}{m^3}$, yaitu sekitar
\qty{35}{km^3}. (2) $100/15.5 \approx 6.5$ tahun.
\end{solution}

\begin{solution}{exo:g12:destruction-of-mountains:12}
$E = kh$ dan $h = h_0 e^{-fkt}$ memberikan $E(t) = kh_0 e^{-fkt}$. Tebal yang
dihilangkan dari $t = 0$ sampai $t \to \infty$ adalah $\int_0^{\infty} kh_0
e^{-fkt}\,\mathrm{d}t = kh_0/(fk) = h_0/f$. Untuk menurunkan permukaan sebesar
$h_0$, erosi harus menghilangkan $h_0/f \approx 6.1\,h_0$, seperti yang
dinyatakan proposisi tentang pantulan: setiap meter yang terkikis hanya
menurunkan permukaan $f$ meter.
\end{solution}

\begin{solution}{exo:g12:destruction-of-mountains:13}
(1) $10 \times 1000 = \qty{10000}{mm}$, yaitu \qty{10}{m}. (2) Menghilangkan
beban es setebal $t$ sama dengan menghilangkan massa $\rho_{\text{es}}\,t$ per
satuan luas; kolom itu naik sampai ia memindahkan mantel sebanyak itu lebih sedikit:
$2000 \times 900/3380 \approx \qty{530}{m}$. (3) Daratan naik makin sedikit
ketika mendekati keseimbangan barunya: yang mendorong kenaikan itu adalah bagian
yang belum pulih.
\end{solution}

\begin{solution}{exo:g12:destruction-of-mountains:14}
(1) Sebelum 1900: sekitar \qty{56}{m/Ma} (seperti dalam contoh). Pada
1987--2006: $145/0.85 \approx 171$ juta ton, $\num{1.71e11}/2830
\approx \qty{6.0e7}{m^3}$, dibagi $\qty{2.98e12}{m^2}$: sekitar
\qty{20}{m/Ma}. (2) Yang lebih awal: muatan yang kemudian tertahan oleh
bendungan dan bangunan lain. (3) Tidak: penurunannya $400 - 145 = 255$ juta ton
setahun, dan bendungan menjelaskan 100 sampai 150 di antaranya, sekitar separuhnya.
\end{solution}

\begin{solution}{exo:g12:destruction-of-mountains:15}
(1) Setidaknya \qty{6}{km}. (2) Paling lama $310 - 252 = 58$ juta tahun, sehingga
setidaknya $6000/58 \approx \qty{100}{m/Ma}$. (3) Laju itu terletak antara median
(\qty{54}{m/Ma}) dan rata-rata (\qty{218}{m/Ma}) cekungan: laju yang biasa
untuk sebuah sabuk pegunungan.
\end{solution}

\begin{solution}{pb:g12:destruction-of-mountains:1}
\textbf{1.} $E = 0.209 \times 2000 \approx \qty{418}{m/Ma}$.

\textbf{2.} \qty{0.418}{mm} setahun.

\textbf{3.} Lebih cepat daripada keduanya: pegunungan itu curam dan muda,
secepat cekungan-cekungan aktif seismik, bahkan lebih.

\textbf{4.} Kemiringan rata-ratanya.

\textbf{5.} Lapisan setebal $\qty{0.418e-3}{m}$ di atas $\qty{1e6}{m^2}$ adalah
\qty{418}{m^3}, $418 \times 2830 \approx \qty{1.18e6}{kg}$: sekitar 1200
ton.

\textbf{6.} $2000/418 \approx 4.8$ juta tahun.

\textbf{7.} $h(t) = 2000\,e^{-0.209\,t}$ ($t$ dalam juta tahun).

\textbf{8.} $1/0.209 \approx 4.8$ juta tahun.

\textbf{9.} $\ln 2/0.209 \approx 3.3$ juta tahun.

\textbf{10.} $e^{-0.209 \times 11} = e^{-2.30} \approx 0.100$.

\textbf{11.} $1 - 2830/3380 \approx 0.163$.

\textbf{12.} Permukaan turun $100 \times 0.163 \approx \qty{16}{m}$;
\omterm{def:g12:continental-crust:root}{akarnya} naik sekitar \qty{84}{m}.

\textbf{13.} $h' = -0.163 \times 0.209\,h \approx -0.0341\,h$, sehingga $h(t) =
2000\,e^{-0.0341\,t}$.

\textbf{14.} Batuan setebal $(2000 - 200)/0.163 \approx \qty{11000}{m}$.

\textbf{15.} Batuan yang semula berada sekitar \qty{11}{km} di bawah permukaan:
batuan metamorf seperti sekis dan gneis, serta granit.

\textbf{16.} $1/0.0341 \approx 29$ juta tahun.

\textbf{17.} $29 \times \ln 10 \approx 67$ juta tahun.

\textbf{18.} $2000\,e^{-0.0341 \times 200} \approx 2000 \times 0.0011 \approx
\qty{2}{m}$.

\textbf{19.} $2000\,e^{-0.0341 \times 250} \approx 2000 \times 0.0002 \approx
\qty{0.4}{m}$: dataran yang rata. Puncak setinggi \qty{2037}{m} berarti bahwa
model itu sendiri tidak cocok: wilayah itu pasti telah terangkat kembali sejak
saat itu, atau terkikis lebih lambat daripada yang diandaikan model.

\textbf{20.} Ketinggian rata-ratanya menjadi separuh dalam $29 \times \ln 2 \approx$
\textbf{20 juta tahun}.
\end{solution}
